\documentclass[../main.tex]{subfiles}
\begin{document}
\label{chapter:appendixC}
Phụ lục này trích dẫn một số đoạn mã nguồn cốt lõi của hệ thống nhằm minh hoạ cách hiện thực
các cơ chế đã trình bày ở các chương trước. Mã được rút gọn (lược bỏ phần khai báo kiểu và
chú thích) để tập trung vào phần logic chính.

\section{Ký và tạo nội dung mã QR}

Đoạn mã trong Danh sách mã nguồn~\ref{lst:c-crypto} (tệp \texttt{src/utils/crypto.ts}) thực hiện việc ký
nội dung mã QR bằng HMAC-SHA256 và tạo chuỗi JSON cuối cùng để mã hoá thành mã QR.

\begin{lstlisting}[caption={crypto.ts --- ký và tạo nội dung mã QR},captionpos=b,label=lst:c-crypto]
import CryptoJS from "crypto-js";

export function signPayload(data, secret) {
  const message =
    `${data.type}:${data.sessionId}:${data.userId}:${data.timestamp}:${data.nonce}`;
  return CryptoJS.HmacSHA256(message, secret).toString();
}

export function createQRContent(type, sessionId, userId, secret) {
  const timestamp = Date.now();
  const nonce = CryptoJS.lib.WordArray.random(16).toString();
  const signature =
    signPayload({ type, sessionId, userId, timestamp, nonce }, secret);
  return JSON.stringify({ type, sessionId, userId, timestamp, nonce, signature });
}
\end{lstlisting}

\section{Sinh mã QR xoay theo thời gian}

Danh sách mã nguồn~\ref{lst:c-qrgen} (tệp \texttt{useQRGenerator.ts}) sinh lại mã QR sau mỗi
30 giây dựa trên một mốc đồng hồ chung (anchor), nhờ đó mã trên Mini App và trên máy chiếu
luôn đồng bộ.

\begin{lstlisting}[caption={useQRGenerator --- mã QR xoay 30 giây có đồng bộ},captionpos=b,label=lst:c-qrgen]
export function useQRGenerator(options) {
  const [qrDataURL, setQrDataURL] = useState("");
  const lastPeriodRef = useRef(-1);
  const refreshInterval = options?.refreshIntervalMs || 30000;

  useEffect(() => {
    if (!options) return;
    const tick = async () => {
      const { type, sessionId, userId, secret, anchor } = options;
      // clock anchor keeps Mini App and projector in sync
      const useAnchor = anchor || Date.now();
      const elapsed = Math.max(0, Date.now() - useAnchor);
      const periodIndex = Math.floor(elapsed / refreshInterval);
      // regenerate QR only when entering a new 30s period
      if (periodIndex !== lastPeriodRef.current) {
        lastPeriodRef.current = periodIndex;
        const content = type === "teacher"
          ? createTeacherQR(sessionId, userId, secret)
          : createPeerQR(sessionId, userId, secret);
        setQrDataURL(await generateQRDataURL(content));
      }
    };
    tick();
    const id = setInterval(tick, 500);
    return () => clearInterval(id);
  }, [options?.sessionId, options?.secret, refreshInterval]);

  return { qrDataURL };
}
\end{lstlisting}

\section{Bảo mật phía máy chủ: chống phát lại và giới hạn tần suất}

Danh sách mã nguồn~\ref{lst:c-server} (tệp \texttt{attendance.service.ts} phía máy chủ) cho thấy hàm
Cloud Function \texttt{scanTeacher} kiểm tra giới hạn tần suất, theo dõi nonce để chống phát
lại và xác minh chữ ký HMAC ở phía máy chủ.

\begin{lstlisting}[caption={Cloud Function --- chống phát lại và giới hạn tần suất},captionpos=b,label=lst:c-server]
// In-memory nonce tracking to prevent QR replay (TTL 120s > QR expiry 90s)
const usedNonces = new Map();
const NONCE_TTL_MS = 120_000;

function checkAndRecordNonce(nonce) {
  const now = Date.now();
  for (const [key, expiry] of usedNonces) {
    if (now > expiry) usedNonces.delete(key);
  }
  if (usedNonces.has(nonce)) return false;
  usedNonces.set(nonce, now + NONCE_TTL_MS);
  return true;
}

export const scanTeacher = onCall(requireAuth(
  async (data, ctx, userId) => {
  if (!checkRateLimit(userId, 10, 60_000)) {
    throw new HttpsError("resource-exhausted", "Too many requests");
  }
  const { qrPayload, sessionId } = data;
  // load session secret, then verify HMAC + timestamp + nonce
  if (!verifyHMAC(qrPayload, secret)) { /* reject: bad signature */ }
  if (!checkAndRecordNonce(qrPayload.nonce)) { /* reject: replay */ }
}));
\end{lstlisting}

\section{Tính điểm tin cậy}

Danh sách mã nguồn~\ref{lst:c-trust} (tệp \texttt{src/types/index.ts}) là hàm tính điểm tin cậy, hiện
thực trực tiếp quy tắc đã trình bày ở Thuật toán~\ref{algo:trust}.

\begin{lstlisting}[caption={computeTrustPolicy --- tính điểm tin cậy},captionpos=b,label=lst:c-trust]
export function computeTrustPolicy(peerCount, faceVerification, config) {
  const faceReq = config?.faceRequired !== false;
  const peerReq = config?.peerRequired !== false;

  const faceOk = faceVerification?.matched === true
    && (faceVerification.confidence || 0) >= 0.7;
  const faceSkipped = faceVerification?.skipped === true;
  const faceAttempted = !!faceVerification && !faceSkipped;

  const facePass = !faceReq || faceOk || faceSkipped || !faceAttempted;
  const peerPass = !peerReq || peerCount >= 1;

  if (facePass && peerPass) return "present";
  if (facePass || peerPass) return "review";
  return "absent";
}
\end{lstlisting}

\section{Quy tắc bảo mật Firestore}

Danh sách mã nguồn~\ref{lst:c-rules} (tệp \texttt{firestore.rules}) minh hoạ các hàm trợ giúp ngăn người
dùng tự nâng quyền và quy tắc áp dụng cho collection \texttt{users}.

\begin{lstlisting}[caption={firestore.rules --- chống tự nâng quyền},captionpos=b,label=lst:c-rules]
function notChangingRole() {
  return !('role' in request.resource.data) ||
         !('role' in resource.data) ||
         request.resource.data.role == resource.data.role;
}

function newDocSafeRole() {
  // new user: only role "student" is accepted
  return !('role' in request.resource.data) ||
         request.resource.data.role == 'student';
}

match /users/{userId} {
  allow read:   if true;
  allow create: if newDocSafeRole();   // no self-elevation
  allow update: if notChangingRole();  // cannot change own role
  allow delete: if false;
}
\end{lstlisting}

\end{document}
